% ==============================================================================
% Projeto de Sistema - Nome do Aluno
% Capítulo 2 - Plataforma de Desenvolvimento
% ==============================================================================
\chapter{Plataforma de Desenvolvimento}
\label{sec-plataforma}
\vspace{-1cm}



%=======================================================================================================
%			Tabela de Plataforma de Desenvolvimento e Tecnologias Utilizadas
%=======================================================================================================

Na Tabela~\ref{tabela-plataforma} são listadas as tecnologias utilizadas no desenvolvimento da ferramenta, bem como o propósito de sua utilização.

\begin{footnotesize}
\begin{longtable}{|p{2.5cm}|c|p{5cm}|p{5.5cm}|}
	\caption{Plataforma de Desenvolvimento e Tecnologias Utilizadas.}	
	\label{tabela-plataforma}\\\hline

	\rowcolor{lightgray}
	\textbf{Tecnologia} & \textbf{Versão} & \textbf{Descrição} & \textbf{Propósito} \\\hline 
	\endfirsthead
	\hline
	\rowcolor{lightgray}
	\textbf{Tecnologia} & \textbf{Versão} & \textbf{Descrição} & \textbf{Propósito} \\\hline 
	\endhead
		
	Next.js & 16.3 & Framework React para desenvolvimento web full-stack, com suporte a renderização no servidor (SSR/SSG) e rotas de API integradas & Base do desenvolvimento da aplicação, unificando front-end e back-end em um único projeto \\ \hline
	TypeScript & 7.0 & Superset do JavaScript que adiciona tipagem estática ao código & Garantir maior segurança e manutenibilidade do código, reduzindo erros em tempo de execução \\ \hline
	PostgreSQL & 18.0 & Sistema de gerenciamento de banco de dados relacional (SGBD) open-source & Armazenamento e persistência dos dados da aplicação \\ \hline
	Drizzle ORM & 0.45 & ORM (\textit{Object-Relational Mapper}) leve e \textit{type-safe} para TypeScript & Mediar a comunicação entre a aplicação e o banco PostgreSQL de forma tipada, sem necessidade de SQL manual \\ \hline
	Better Auth & --- & Biblioteca de autenticação para TypeScript que executa dentro da própria aplicação e guarda usuários, sessões e credenciais no banco de dados, com adaptador para o Drizzle ORM & Cadastro e login por e-mail e senha, gerenciamento de sessões, armazenamento seguro das senhas e proteção das telas e \textit{Server Actions} conforme o perfil do usuário (Tutor, Passeador ou Atendente) \\ \hline
	Tailwind CSS & 4.3 & Framework CSS \textit{utility-first} para estilização de interfaces & Estilização rápida e consistente dos componentes visuais do sistema \\ \hline
	shadcn/ui & --- & Coleção de componentes de interface reutilizáveis sobre Radix UI e Tailwind, instalada via CLI & Acelerar o desenvolvimento do front-end com componentes prontos e customizáveis \\ \hline

\end{longtable}
\end{footnotesize}


%=======================================================================================================
%			Tabela de Softwares de Apoio ao Desenvolvimento do Projeto
%=======================================================================================================

Na Tabela~\ref{tabela-software} vemos os softwares que apoiaram o desenvolvimento de documentos e também do código fonte.

\begin{footnotesize}
\begin{longtable}{|p{2.5cm}|c|p{5cm}|p{5.5cm}|}
	\caption{Softwares de Apoio ao Desenvolvimento do Projeto}	
	\label{tabela-software}\\\hline
	
	\rowcolor{lightgray}
	\textbf{Tecnologia} & \textbf{Versão} & \textbf{Descrição} & \textbf{Propósito} \\\hline 
	\endfirsthead
	\hline
	\rowcolor{lightgray}
	\textbf{Tecnologia} & \textbf{Versão} & \textbf{Descrição} & \textbf{Propósito} \\\hline 
	\endhead
		
	WebStorm & 2026.2 & IDE da JetBrains voltada para desenvolvimento JavaScript/TypeScript, com suporte nativo a frameworks como Next.js e React & Escrita, depuração e navegação no código-fonte da aplicação \\ \hline
	Git & 2.55 & Sistema de controle de versão distribuído & Versionamento do código-fonte e rastreamento do histórico de alterações \\ \hline
	GitHub & --- & Plataforma de hospedagem de repositórios Git, com suporte a \textit{pull requests}, \textit{issues} e integração contínua & Hospedagem do repositório remoto e colaboração entre os membros da equipe \\ \hline
	Overleaf (LaTeX) & --- & Plataforma online de edição colaborativa de documentos em LaTeX & Elaboração e formatação da documentação do projeto (como este documento) \\ \hline
	Trello & --- & Ferramenta de gerenciamento de projetos baseada em quadros \textit{Kanban} & Organização, distribuição e acompanhamento das tarefas da equipe \\ \hline
\end{longtable}
\end{footnotesize}